\section*{Введение}
\addcontentsline{toc}{section}{Введение}

\textbf{Актуальность темы.} Разговорная практика является ключевым, но наиболее
труднодоступным компонентом изучения иностранного языка. Аудиторных занятий в
университете недостаточно для формирования навыка свободной речи, а самостоятельный
поиск собеседника или языкового клуба затруднён из-за рассредоточенности информации
по чатам и социальным сетям. Разработка специализированной веб-платформы,
объединяющей каталог языковых клубов и инструменты записи на встречи, позволяет
устранить этот барьер.

\textbf{Степень разработанности проблемы.} Существующие решения -- приложения для
поиска собеседника (HelloTalk, Tandem) и универсальные платформы мероприятий
(Meetup) -- решают задачу частично: первые не поддерживают групповой формат, вторые
не имеют языковой специализации. Специализированного веб-сервиса, сочетающего
клубный формат, уровни владения языком и инструменты записи на встречи, на рынке
не представлено.

\textbf{Объект исследования} -- процессы организации разговорной практики
английского языка в формате языковых клубов.

\textbf{Предмет исследования} -- программные средства для автоматизации поиска
языковых клубов и управления записью участников на встречи в рамках веб-платформы.

\textbf{Целью данной работы} является проектирование и разработка веб-приложения
для поиска языковых клубов и записи на разговорные встречи по английскому языку
с разграничением доступа по ролям пользователей.
